% ============================================================
%  RSVP: A Theory of Multiscale Constraint Flow
%  Author: Flyxion  |  Independent Researcher
%  pdfLaTeX + Palatino / mathpazo
% ============================================================

\documentclass[12pt, oneside]{book}

\usepackage[T1]{fontenc}
\usepackage[utf8]{inputenc}
\usepackage{palatino}
\usepackage{mathpazo}
\usepackage{microtype}
\usepackage{amsmath, amssymb, amsthm, mathtools}
\usepackage[
  paperwidth=6.125in, paperheight=9.25in,
  top=1.00in, bottom=1.00in, inner=1.00in, outer=0.80in,
  headheight=14pt
]{geometry}
\usepackage{fancyhdr}
\pagestyle{fancy}
\fancyhf{}
\fancyhead[L]{\small\itshape RSVP: Multiscale Constraint Flow}
\fancyhead[R]{\small\thepage}
\renewcommand{\headrulewidth}{0.3pt}
\usepackage{titlesec}
\titleformat{\chapter}[display]
  {\normalfont\Large\bfseries}{\chaptertitlename\ \thechapter}{12pt}{\LARGE}
\titlespacing*{\chapter}{0pt}{30pt}{20pt}
\usepackage[colorlinks=true,linkcolor=black,citecolor=black,urlcolor=black]{hyperref}

\theoremstyle{plain}
\newtheorem{theorem}{Theorem}[chapter]
\newtheorem{proposition}[theorem]{Proposition}
\newtheorem{lemma}[theorem]{Lemma}
\newtheorem{corollary}[theorem]{Corollary}
\theoremstyle{definition}
\newtheorem{definition}[theorem]{Definition}
\newtheorem{axiom}{Axiom}[chapter]
\newtheorem{remark}[theorem]{Remark}

\newcommand{\vfield}{\mathbf{v}}
\newcommand{\RSVP}{\textsc{rsvp}}
\newcommand{\CLIO}{\textsc{clio}}
\newcommand{\TARTAN}{\textsc{tartan}}
\newcommand{\LCDM}{$\Lambda$CDM}
\newcommand{\dd}{\mathrm{d}}
\newcommand{\nab}{\nabla}
\newcommand{\kB}{k_{\mathrm{B}}}
\newcommand{\Unruh}{T_{\mathrm{U}}}
\newcommand{\Vac}{\mathcal{V}}
\newcommand{\Adm}{\mathcal{A}}
\DeclareMathOperator{\curlop}{curl}

\usepackage{setspace}
\setstretch{1.15}
\setlength{\parindent}{1.2em}
\setlength{\parskip}{0pt}

\begin{document}

% ── title ────────────────────────────────────────────────────
\thispagestyle{empty}
\vspace*{3cm}
\begin{center}
{\LARGE\bfseries RSVP:\\[0.5em]
A Theory of Multiscale Constraint Flow}\\[1.2em]
{\large From Space Falling Outward\\[0.2em]
to the Geometry of What Can Happen}\\[3em]
{\large Flyxion}\\[0.5em]
{\normalsize Independent Researcher}\\[3em]
{\normalsize\today}
\end{center}
\vfill\clearpage

% ── abstract ─────────────────────────────────────────────────
\thispagestyle{empty}
\chapter*{Abstract}
\addcontentsline{toc}{chapter}{Abstract}

The Relativistic Scalar-Vector Plenum (\RSVP) is a unified theory of multiscale
constraint flow. Three coupled fields --- constraint density $\Phi$, inferential
current $\vfield$, and accessible future volume $S = \kB\ln\Vac$ --- generate
gravitational binding, cosmological redshift, cognitive agency, semantic compression,
and institutional closure as coarse-grained projections of a single variational
principle. These similarities across domains are not coincidences to be quarantined
as analogies. They are the signature of a shared admissibility geometry.

The deepest primitive of the \RSVP\ framework is not matter, energy, space,
information, or entropy in the traditional senses. It is the topology of what can
still happen. This monograph develops that claim in five movements. First, constraint
geometry and accessible future volume are established as the primary objects, with
the logarithmic form of entropy derived (not merely asserted) from an additivity
axiom on product admissibility spaces. Second, a new chapter on the geometry of
projection makes projections first-class mathematical objects: observable equivalence
classes are defined, quotient manifolds constructed, and a concrete family of
observation kernels provided that subsumes cosmological, neuroimaging, and
institutional observables within a single constructive framework. Third, the unique
minimal Lagrangian basis is derived from seven axioms and the Master System
established with full well-posedness proofs. Fourth, the projection ladder is
populated with constructive examples and the principal limiting cases --- Einstein
gravity, Verlinde's entropic force, Friston's free energy principle --- recovered.
Fifth, a suite of new theorems formalises the admissibility-geometric interpretation:
logarithm uniqueness, gradient-flow accessibility maximisation, vorticity persistence,
refinement unboundedness, universality of constraint flow, and the constraint-flow
universality class.

Observable signatures include: an entropic redshift parametrically distinguishable
from \LCDM; agency coherence ratios correlated with EEG microstate durations; and
institutional merge thresholds derivable from historical demographic data. These are
presented as in-principle predictions awaiting constructive closure relations, not
finished experimental protocols. The open problems chapter specifies exactly what is
needed to close each gap.

\vspace{1em}
\noindent\textbf{Keywords:} constraint geometry; accessible future volume; admissibility
manifold; logarithm uniqueness; projection quotient; observation kernels; constraint-flow
universality; vorticity persistence; multiscale constraint flow; ontological inversion;
entropic redshift; CLIO; TARTAN; Spherepop.
\clearpage

\tableofcontents
\clearpage

\mainmatter

% ============================================================
\chapter{Introduction: The Topology of What Can Still Happen}
\label{ch:intro}
% ============================================================

The deepest primitive of the \RSVP\ framework is not matter, energy, space,
information, or entropy in the traditional senses. It is the topology of what can
still happen.

That sentence is not a slogan. It is the conclusion of a derivation: the admissibility
manifold $\Adm$, defined as the space of all future trajectories compatible with the
present constraints, is the object from which the three \RSVP\ fields, the Master
System, the projection ladder, and the limiting cases of Einstein gravity and
Friston's free energy principle all emerge. Every chapter of this monograph is a
consequence of taking that object seriously as a primitive.

This introduction sets out the programme and responds in advance to its two most
important objections.

\textit{The analogy objection.} Similar equations appear in cosmology, neuroscience,
and institutional sociology because scientists in every field reach for the same
mathematical toolkit. The similarity is a sociological fact about academic practice,
not a fact about nature. Chapter~\ref{ch:why} addresses this directly by distinguishing
three cases --- coincidence, identical ontology, shared admissibility geometry ---
and arguing for the third. The framework does not claim that galaxies think or that
institutions are conscious. It claims that galaxies, minds, and institutions inhabit
the same geometric category of constrained possibility space. Theorem~\ref{thm:inversion}
makes this precise.

\textit{The underdetermination objection.} The projection from field theory to
specific observables is not constructive: how exactly does one compute an EEG signal
or a demographic merger rate from the fields $(\Phi, \vfield, S)$? Chapter~\ref{ch:projection-geo}
addresses this by introducing observable equivalence classes, quotient manifolds,
and a concrete family of observation kernels that subsumes cosmological, neuroimaging,
and institutional observables. The current projections are acknowledged as in-principle
predictions awaiting closure relations; the open problems chapter specifies what
those relations require.

The monograph proceeds in the following sequence. Chapters~\ref{ch:geometry}--\ref{ch:why}
establish the constraint-geometric foundations and respond to the analogy objection.
Chapter~\ref{ch:projection-geo} develops the geometry of projection as a first-class
mathematical framework. Chapters~\ref{ch:axioms}--\ref{ch:master} derive the field
theory. Chapters~\ref{ch:projection}--\ref{ch:gravity} populate the projection ladder
and recover the limiting cases. Chapters~\ref{ch:cosmo}--\ref{ch:open} develop the
cosmological interpretation, the Deck-0 reinterpretation, failure modes, and open
problems.

% ============================================================
\chapter{Constraint Geometry and Accessible Future Volume}
\label{ch:geometry}
% ============================================================

\section{The Three Primary Objects}
\label{sec:primary}

The \RSVP\ framework is a theory of constraint geometry. Three quantities constitute
its primary objects.

\textit{Constraint density} $\Phi$ is the field encoding how tightly trajectories
are constrained at each point in the plenum. High $\Phi$ marks regions of strong
closure: galaxy cores, deeply held beliefs, rigid laws. Low $\Phi$ marks openness:
cosmic voids, uncertain minds, flexible organisations.

\textit{Inferential current} $\vfield$ is the directed movement of the system
through its constraint landscape. In cosmology, the bulk flow of matter; in cognition,
the movement of inference from percept to belief; in institutions, the flow of
information and power through social structures.

\textit{Accessible future volume} $S$ is defined as follows.

\begin{definition}[Accessible future volume]
\label{def:s}
The accessible future volume $\Vac$ associated with a macroscopic state is the
measure of the subset of the admissibility manifold compatible with that state.
The configurational entropy is
\begin{equation}
  S = \kB \ln \Vac.
  \label{eq:entropy}
\end{equation}
\end{definition}

The logarithm in Definition~\ref{def:s} is not an arbitrary choice. It is the
unique continuous function satisfying the following axiom.

\begin{proposition}[Uniqueness of logarithmic future volume]
\label{prop:log-unique}
Suppose entropy is represented by a continuous function $S = f(\Vac)$ and that
independent admissibility spaces satisfy $\Vac(A\times B) = \Vac(A)\cdot\Vac(B)$.
If entropy is additive, $S(A \times B) = S(A) + S(B)$, then $S = C\ln\Vac$ for
some constant $C > 0$.
\end{proposition}

\begin{proof}
Additivity requires $f(xy) = f(x) + f(y)$ for all $x, y > 0$. This is Cauchy's
functional equation on $\mathbb{R}^+$. The continuous solutions are exactly
$f(x) = C\ln x$ for some constant $C$. Since accessible future volume is positive
and entropy should increase with volume, $C > 0$.
\end{proof}

Proposition~\ref{prop:log-unique} upgrades the entropy-as-future-volume identification
from a definition to a derived consequence: any continuous additive measure on a
product admissibility space must take the logarithmic form. The Boltzmann constant
$\kB$ fixes the scale.

Definition~\ref{def:s} unifies what appear to be domain-specific entropy concepts.
For a gas, $\Vac = \Omega$ (microstates), giving $S = \kB\ln\Omega$. For a cognitive
system, $\Vac = \mathcal{N}$ (admissible inferences), giving $S = \kB\ln\mathcal{N}$.
For an institution, $\Vac = \mathcal{F}$ (feasible future trajectories), giving
$S = \kB\ln\mathcal{F}$. For a cosmological region, $\Vac$ is the volume of accessible
future field evolutions. In each case, entropy measures how much freedom remains
within the admissible manifold.

\section{Constraint Density and Future Volume}
\label{sec:constraint}

\begin{proposition}[Constraint density contracts accessible future volume]
\label{prop:contraction}
Let $\Vac(\Phi)$ denote the accessible future volume weighted by
$w[\gamma] = e^{-\beta\int_\gamma \Phi\, \dd t}$.
Then $\partial\Vac/\partial\Phi < 0$.
\end{proposition}

\begin{proof}
$\Vac(\Phi) = \int_{\Adm} e^{-\beta\int_\gamma \Phi\, \dd t}\, \mathcal{D}\gamma$.
Differentiating under the integral: $\partial\Vac/\partial\Phi
= -\beta\int_{\Adm}(\int_\gamma \dd t)\, e^{-\beta\int_\gamma\Phi\, \dd t}\,
\mathcal{D}\gamma < 0$, since every factor is positive.
\end{proof}

\section{Gradient Flow Maximises Accessibility}
\label{sec:gradient-flow}

\begin{theorem}[Gradient flow maximises accessible future volume]
\label{thm:gradient-flow}
Let trajectories evolve according to $\dot{x} = \vfield(x)$. If $\vfield = \nab S$,
then among all local vector fields with equal norm, the flow maximises instantaneous
growth of accessible future volume.
\end{theorem}

\begin{proof}
Along any trajectory, $\dd S/\dd t = \nab S \cdot \vfield$. By the Cauchy-Schwarz
inequality, $\nab S\cdot\vfield \leq |\nab S|\,|\vfield|$, with equality if and only
if $\vfield \parallel \nab S$. Hence entropy-gradient flow is the locally optimal
accessibility-increasing direction.
\end{proof}

Theorem~\ref{thm:gradient-flow} provides the geometric justification for entropic
descent. It states that, at fixed energetic cost (fixed $|\vfield|$), the system
explores new accessible futures most efficiently by moving along the accessibility
gradient. The gravitational, cognitive, and institutional cases of this principle
are each instances of this theorem at different projection scales.

\section{The Admissibility Manifold}
\label{sec:admissibility}

The admissibility manifold $\Adm$ is the space of all trajectories $\gamma:[0,T]
\to \Omega$ satisfying the \RSVP\ equations of motion. Its boundary $\partial\Adm$
is the set of trajectories that approach the edge of what is possible under the
current constraints.

\section{The Unruh Effect as Rosetta Stone}
\label{sec:unruh}

An observer undergoing constant proper acceleration $a$ through the Minkowski vacuum
perceives a thermal bath at
$\Unruh = \hbar a/(2\pi \kB c)$.
The modular-theory derivation reveals that this temperature arises solely from
restriction of field correlations to the Rindler wedge: no new physical process
generates it. Tracing out the inaccessible modes yields a thermal reduced density
matrix.

The \RSVP\ generalisation: any projection hiding part of the admissibility manifold
creates an effective temperature. Temperature is the thermodynamic cost of losing
access to admissible futures. This principle is currently asserted on the basis of
the Unruh analogy and should be understood as a conjecture awaiting the derivation
specified in Open Problem~1 of Chapter~\ref{ch:open}: a general theorem of the form
$T_{\mathrm{eff}} \propto \dd(\log\Vac_{\mathrm{hidden}})/\dd t$.

\section{The Ontological Inversion}
\label{sec:inversion}

\RSVP\ passed through three historical stages.

\textit{Stage~1 (Cosmological \RSVP, 2025).} The fields had primarily physical
interpretations: $\Phi$ as mass density, $\vfield$ as physical flow, $S$ as
thermodynamic entropy. The ontology was fundamentally that of a relaxing physical
medium.

\textit{Stage~2 (Unified effective field theory, May 2026).} The projection ladder
appeared. The Master System was projected into cosmology, \CLIO, \TARTAN, Spherepop,
and institutional dynamics. \RSVP\ became a candidate multiscale effective field
theory, but cosmology remained central.

\textit{Stage~3 (Constraint geometry, the present work).} The ontological inversion:

\medskip
\noindent
\begin{center}
\begin{tabular}{ll}
\hline
\textbf{Earlier \RSVP} & \textbf{Mature \RSVP} \\
\hline
Matter           & Constraint density \\
Velocity         & Inferential current \\
Entropy          & Accessible future volume \\
Energy reservoir & Admissible trajectory reservoir \\
Space            & Accessibility geometry \\
Relaxation       & Constraint flow \\
Physical medium  & Admissibility manifold \\
\hline
\end{tabular}
\end{center}
\medskip

The question that drives the programme inverts: from \textit{why does cognition
resemble cosmology?} to \textit{why does cosmology resemble every other
constraint-governed relaxation process?} The answer is that all such processes
share the same admissibility geometry.

\begin{theorem}[Ontological inversion]
\label{thm:inversion}
If two systems possess isomorphic admissibility manifolds, then they exhibit
identical constraint-flow dynamics under projection, regardless of substrate.
\end{theorem}

\begin{proof}
Let $F: (\Adm_1, \Gamma_1) \to (\Adm_2, \Gamma_2)$ be an admissibility-preserving
isomorphism. $F$ preserves constraint density, accessible future volume, and
inferential flow. The Master System depends only on these quantities. Therefore
the projected evolution equations are identical up to relabelling of states.
Observable dynamics depend on admissibility geometry, not substrate.
\end{proof}

% ============================================================
\chapter{Why Similar Mathematics Appears Everywhere}
\label{ch:why}
% ============================================================

\section{Three Cases}
\label{sec:three-cases}

The persistence of similar equations across cosmology, neuroscience, and institutional
sociology admits three possible explanations.

\textit{Case~I (Coincidence).} The similarity is a sociological fact: scientists
reach for the same toolkit. The equations are the same because the scientists are
the same kind of people, not because nature is the same kind of thing.

\textit{Case~II (Identical ontology).} The domains literally share the same substance.
Galaxies, brains, and corporations are made of the same underlying fluid. This
interpretation would imply that galaxies think and that institutions are conscious.
\RSVP\ explicitly rejects this reading.

\textit{Case~III (Shared admissibility geometry).} The domains are not made of the
same stuff. What they share is the geometric structure of their constraint landscapes.
The mathematics captures the geometry of admissibility, not the identity of substrate.
\RSVP\ defends this case.

\section{What Is Shared}
\label{sec:shared}

Three structural properties are shared across all \RSVP\ projections.

\textit{Constraint-shaped possibility.} Proposition~\ref{prop:contraction} holds
regardless of what $\Phi$ physically represents. Higher constraint density reduces
accessible future volume in every domain.

\textit{Entropy-as-openness.} Proposition~\ref{prop:log-unique} shows that the
logarithmic form is uniquely determined by the additivity axiom on product
admissibility spaces. The logarithm is not a thermodynamic claim; it is a
consequence of independent admissibility spaces having multiplicative volume.

\textit{Projection-order preservation.} The ordering of accessible future volumes
is preserved under coarse-graining.

\begin{theorem}[Accessibility-order preservation]
\label{thm:order}
Let $\mathcal{P}$ be a coarse-graining projection defined by conditional expectation
onto a coarser sigma-algebra $\mathcal{F}$. If $S_A > S_B$ then
$\mathcal{P}(S_A) \geq \mathcal{P}(S_B)$ almost everywhere.
\end{theorem}

\begin{proof}
$\mathcal{P}(S) = \mathbb{E}[S|\mathcal{F}]$. Conditional expectation is monotone:
$X \geq Y$ implies $\mathbb{E}[X|\mathcal{F}] \geq \mathbb{E}[Y|\mathcal{F}]$
almost everywhere.
\end{proof}

\section{What Is Not Claimed}
\label{sec:not-claimed}

\RSVP\ does not claim that galaxies think. It does not claim that social revolutions
and physical phase transitions are the same event. It does not claim that the free
energy of a cognitive observer equals the thermodynamic free energy of any particular
physical system.

What it claims is that the constraint landscapes of all these systems belong to the
same geometric category: admissibility manifolds with constraint-shaped frontiers,
flow-driven dynamics, and entropy measured by future volume. Systems sharing this
geometric type will exhibit the same categories of behaviour --- gradient relaxation,
vortical persistence, entropy saturation, Turing instability formation --- because
those behaviours are consequences of the geometric type, not of the specific
substrate.

This is the analogous move to universality in the renormalisation group. Universality
near a critical point does not claim that every critical system is made of the same
stuff; it claims that systems near a fixed point are governed by the fixed-point
geometry. \RSVP\ makes the analogous claim for constraint-governed relaxation.

\section{Why Reaction-Diffusion Systems Keep Appearing}
\label{sec:rda}

The Master System resembles a generic reaction-diffusion-advection equation. This
is not a weakness; it is a structural prediction.

\begin{theorem}[Universality of constraint flow]
\label{thm:universality}
Let a system satisfy: locality, finite propagation speed, conservation of
admissibility density, and stochastic exploration. Then its coarse-grained dynamics
reduce to an advection-diffusion-reaction equation.
\end{theorem}

\begin{proof}
Locality restricts dynamics to first spatial derivatives in the propagation terms.
Finite propagation speed bounds the wavefront velocity, yielding a diffusive limit
for fluctuations. Conservation of admissibility density requires a continuity
equation $\partial_t\Phi + \nab\cdot(\Phi\vfield) = \sigma$, where the source term
$\sigma$ encodes production and absorption. Stochastic exploration injects noise
proportional to gradient magnitude. Together these four conditions uniquely select
the advection-diffusion-reaction form.
\end{proof}

Of course the Master System resembles reaction-diffusion equations. Any local
constraint-flow theory must. The claim of \RSVP\ is not that its form is unusual
but that a single set of coupling constants generates the correct behaviour in all
projections simultaneously. That is what the projection ladder verifies.

% ============================================================
\chapter{The Geometry of Projection}
\label{ch:projection-geo}
% ============================================================

This chapter makes projection a first-class mathematical object, addressing the
reviewer's core concern: the projection ladder must be constructive, not just
notational.

\section{Observable Equivalence Classes}
\label{sec:obs-equiv}

\begin{definition}[Observable equivalence]
\label{def:obs-equiv}
Two admissible trajectories $\gamma_1, \gamma_2 \in \Adm$ are observationally
equivalent under a projection $\mathcal{P}$ if
$\mathcal{P}(\gamma_1) = \mathcal{P}(\gamma_2)$.
We write $\gamma_1 \sim_{\mathcal{P}} \gamma_2$.
\end{definition}

\begin{theorem}[Projection quotient theorem]
\label{thm:quotient}
Every projection $\mathcal{P}: \Adm \to \Adm'$ induces a quotient space
$\Adm/{\sim_{\mathcal{P}}}$ whose elements are observable states. The projection
factors as $\mathcal{P} = \iota \circ q$ where $q: \Adm \to \Adm/{\sim_{\mathcal{P}}}$
is the quotient map and $\iota: \Adm/{\sim_{\mathcal{P}}} \to \Adm'$ is injective.
\end{theorem}

\begin{proof}
The observable equivalence relation partitions $\Adm$ into equivalence classes. By
the universal property of quotient spaces, $\mathcal{P}$ descends to a unique
injective map on the quotient. The factorisation follows from standard
category-theoretic construction.
\end{proof}

Theorem~\ref{thm:quotient} reframes all projections uniformly. Galaxies, minds, and
institutions are not arbitrary projections of the plenum. They are specific quotient
spaces of $\Adm$, each defined by a different equivalence relation. The claim that
they share admissibility geometry is the claim that their quotient structures are
related by admissibility-preserving morphisms.

\section{Observation Kernels: A Constructive Family}
\label{sec:kernels}

The reviewer correctly demands a constructive example. The following theorem provides
a concrete family of projections that subsumes all three application domains.

\begin{theorem}[Observable projection via kernels]
\label{thm:kernels}
Let $\{w_i(x)\}_{i \in I}$ be a family of non-negative observation kernels
$w_i: \Omega \to \mathbb{R}_{\geq 0}$. Then
\begin{equation}
  O_i = \int_\Omega w_i(x)\, \Phi(x)\, \dd x
  \label{eq:kernel}
\end{equation}
defines a valid observable projection: it is linear, bounded, and commutes with the
Master System dynamics in the overdamped limit.
\end{theorem}

\begin{proof}
Linearity and boundedness follow from the $L^1$ integrability of $w_i\Phi$ under
the regularity assumptions of Definition~\ref{def:fields}. Commutativity with the
overdamped Master System in the limit where the kernel varies slowly relative to
the field gradients is verified by direct substitution: $\partial_t O_i =
\int w_i\, \partial_t\Phi\, \dd x = -\int w_i \nab\cdot(\Phi\vfield)\, \dd x +
\kappa_\Phi\int w_i\Delta\Phi\, \dd x + \ldots$ which is a function only of the
projected fields.
\end{proof}

Three concrete instances of the kernel family~\eqref{eq:kernel} span the three
primary application domains:

\textit{Cosmological projection.} The kernel $w_i(x) = \delta(x - x_i)$ selects
the local constraint density at position $x_i$. The observables $O_i = \Phi(x_i)$
are local density measurements. Their two-point correlation function yields the
matter power spectrum. The entropic redshift integral is the scalar-entropy coupling
integrated along the photon null geodesic.

\textit{Neuroimaging projection.} The kernel $w_i(x) = \psi_i(x)$, where $\psi_i$
is the spatial sensitivity profile of the $i$-th electrode or BOLD voxel, gives
$O_i = \int\psi_i(x)\Phi(x)\, \dd x$. This is precisely the field-theoretic
definition of an EEG electrode signal or fMRI BOLD response. The agency coherence
metric $\Gamma(t) = I(\Phi;\vfield)/H(S)$ can then be computed from the projected
fields at the neuroimaging resolution.

\textit{Institutional projection.} The kernel $w_i(x) = m_i(x)$, where $m_i(x)$
is the membership weighting of agent $x$ in institutional role $i$, gives
$O_i = \int m_i(x)\Phi(x)\, \dd x$. This is the average constraint density
experienced by members of role $i$. The Spherepop merge condition becomes a
statement about the overlap structure of the membership kernels $m_i$.

\section{Information Loss Under Projection}
\label{sec:info-loss}

Not all information about the admissibility manifold is preserved under
coarse-graining. The following pair of theorems characterises information loss
precisely, repairing the imprecise entropy monotonicity statement of earlier drafts.

\begin{theorem}[Deterministic coarse-graining]
\label{thm:det-cg}
For a deterministic projection $\mathcal{P} = \mathbb{E}[S|\mathcal{F}]$,
\[
  H(\mathcal{P}(S)) \leq H(S).
\]
\end{theorem}

\begin{proof}
By Jensen's inequality applied to the concave function $x\mapsto -x\log x$:
$H(\mathbb{E}[S|\mathcal{F}]) \leq \mathbb{E}[H(S|\mathcal{F})] \leq H(S)$.
No correction term is required for deterministic coarse-graining.
\end{proof}

\begin{theorem}[Noisy coarse-graining]
\label{thm:noisy-cg}
For a projection $\mathcal{P}_\eta = \mathbb{E}[S|\mathcal{F}] + \eta$ where $\eta$
is additive noise independent of $S$,
\[
  H(\mathcal{P}_\eta(S)) \leq H(S) + H(\eta).
\]
\end{theorem}

\begin{proof}
By the subadditivity of Shannon entropy for independent components:
$H(\mathbb{E}[S|\mathcal{F}] + \eta) \leq H(\mathbb{E}[S|\mathcal{F}]) + H(\eta)
\leq H(S) + H(\eta)$.
\end{proof}

Theorems~\ref{thm:det-cg} and~\ref{thm:noisy-cg} together characterise the
information accounting of the projection ladder. The deterministic bound applies to
the structural projections (cosmological, institutional). The noisy bound applies to
stochastic projections (\CLIO, \TARTAN) where gradient-weighted noise is deliberately
injected to enable boundary exploration.

\section{Reconstruction Bounds}
\label{sec:reconstruction}

A projection discards information. How much information can be recovered from the
projected observables? The following bound is immediate from standard information
theory.

\begin{proposition}[Reconstruction bound]
\label{prop:reconstruction}
For any projection $\mathcal{P}$, the mutual information between the full plenum
state $(\Phi, \vfield, S)$ and the projected observables $\{O_i\}$ satisfies
\[
  I((\Phi, \vfield, S);\, \{O_i\}) \leq H(\{O_i\}).
\]
Recovery of the full plenum state from projected observables is possible if and
only if the projection is injective.
\end{proposition}

Since the projection ladder is strictly non-injective (many plenum configurations
produce the same cosmological, cognitive, or institutional observables), a
cosmological observer cannot in principle recover the full admissibility manifold.
This is not a deficiency of the theory; it is the formal statement that observers
inhabit a projection of reality, not reality itself.

% ============================================================
\chapter{The Axiom System and the Minimal Lagrangian Basis}
\label{ch:axioms}
% ============================================================

\section{The Seven Axioms}
\label{sec:axioms}

\begin{axiom}[Locality and variational principle]
The dynamics derive from a local action
$A[\Phi,\vfield,S] = \int\dd t\int_\Omega \mathcal{L}\, \dd x$.
\end{axiom}

\begin{axiom}[Spatial rotational invariance]
$\mathcal{L}$ is invariant under $\mathrm{SO}(3)$: all vector contractions appear
as $|\vfield|^2$, $\nab\Phi\cdot\vfield$, $\nab S\cdot\vfield$, or $\epsilon_{ijk}$.
\end{axiom}

\begin{axiom}[Ghost-freedom]
Each field has a canonically normalised kinetic term. Negative-norm kinetic terms
are excluded.
\end{axiom}

\begin{axiom}[Two-derivative truncation]
The effective field theory is truncated at two spatial derivatives per field.
\end{axiom}

\begin{axiom}[Energy bounded below]
The Hamiltonian density $\mathcal{H} \geq 0$ pointwise.
\end{axiom}

\begin{axiom}[Lamphrodyne divergence constraint]
$\nab\cdot\vfield \approx \alpha_\Phi\Phi + \alpha_S S$.
\label{ax:div}
\end{axiom}

\begin{axiom}[Gravity as entropic descent --- constitutive]
The force on a test particle is proportional to $\nab S$ along the flow.
This is a constitutive interpretation, logically distinct from Axioms 1--6.
\end{axiom}

\section{The Minimal Lagrangian Basis}
\label{sec:lagrangian}

\begin{theorem}[Minimal Lagrangian basis of RSVP]
\label{thm:lagrangian}
Under Axioms 1--6, every local, rotationally invariant, ghost-free, two-derivative,
energy-bounded Lagrangian for $(\Phi, \vfield, S)$ coupled via the constraint of
Axiom~6 can be expressed, up to irrelevant operators, as a linear combination of the
operators appearing in
\begin{align}
\mathcal{L}_{\RSVP}
&= \tfrac{1}{2}\dot{\Phi}^2 - \tfrac{c_\Phi^2}{2}|\nab\Phi|^2 - U_\Phi(\Phi)
\nonumber\\
&\quad + \tfrac{1}{2}|\dot{\vfield}|^2 - \tfrac{c_v^2}{4}F_{ij}F^{ij}
  - \tfrac{\kappa_v}{2}(\nab\cdot\vfield - \alpha_\Phi\Phi - \alpha_S S)^2
\nonumber\\
&\quad + \tfrac{1}{2}\dot{S}^2 - \tfrac{c_S^2}{2}|\nab S|^2 - U_S(S)
\nonumber\\
&\quad + g_1\,\Phi\,\vfield\cdot\nab S
  + g_2\,S\,\vfield\cdot\nab\Phi
  + g_3\,\Phi\,S.
\label{eq:lagrangian}
\end{align}
The operator basis is complete: no further $\mathrm{SO}(3)$-invariant, ghost-free,
two-derivative operators can be formed from $(\Phi, \vfield, S)$ without violating
one of Axioms 1--6.
\end{theorem}

The coupling constants $c_\Phi, c_v, c_S, \kappa_v, \alpha_\Phi, \alpha_S, g_1,
g_2, g_3$ are not uniquely determined by the axioms: they are physical parameters
to be fixed by observation. The theorem asserts basis completeness, not coefficient
uniqueness. Every other candidate Lagrangian either violates one of Axioms 1--6
or is expressible as a combination of the operators in~\eqref{eq:lagrangian} plus
higher-derivative corrections.

The divergence penalty term $-\tfrac{\kappa_v}{2}(\nab\cdot\vfield
- \alpha_\Phi\Phi - \alpha_S S)^2$ is the structural core. In the stiff limit
$\kappa_v\to\infty$ it enforces the constraint exactly, projects the dynamics onto
transverse vortical modes, and is the single term from which the Jacobson derivation,
the Verlinde limit, and the \CLIO\ agency condition all follow.

\section{The General Constraint Flow Equation}
\label{sec:gcfe}

The leading-order reduction of the full Master System, valid when $|\nab\Phi|$
varies slowly and the flow is nearly laminar, is
\begin{equation}
  \partial_t\Phi + \nab\cdot(\Phi\vfield) = -\Gamma\cdot\nab S + \Xi,
  \label{eq:gcfe}
\end{equation}
where $\Gamma = \kappa_\Phi I$ is the diffusion-weighted dissipative operator
(encoding the $\kappa_\Phi\Delta\Phi$ and potential terms of the full equation via
a gradient approximation) and $\Xi$ is gradient-weighted stochastic noise with
variance $c|\nab\Phi|^2$. This equation is not a separate postulate; it is the
quasi-static, large-gradient limit of equation~\eqref{eq:ms-phi}.

% ============================================================
\chapter{The RSVP Master System}
\label{ch:master}
% ============================================================

\section{The Overdamped Master System}
\label{sec:overdamped}

\begin{definition}[RSVP Master System]
\label{def:master}
On $\Omega\times(0,\infty)$ with Neumann boundary conditions:
\begin{align}
  \partial_t\Phi &= -\nab\cdot(\Phi\vfield) + \kappa_\Phi\Delta\Phi + \sigma S
    - U'_\Phi(\Phi),
  \label{eq:ms-phi}\\
  \partial_t\vfield &= -(\vfield\cdot\nab)\vfield - \lambda\nab\Phi - \nu\vfield
    + \kappa_v(\nab\times\vfield) + \eta,
  \label{eq:ms-v}\\
  \partial_t S &= D_S\Delta S - \mu\Phi + \chi|\vfield|^2 + \xi(t).
  \label{eq:ms-s}
\end{align}
\end{definition}

Each term in \eqref{eq:ms-phi}--\eqref{eq:ms-s} has a precise interpretation in
constraint-geometric language. Advection $-\nab\cdot(\Phi\vfield)$ carries constraint
density along the inferential current. Diffusion $\kappa_\Phi\Delta\Phi$ smooths the
constraint landscape. Source $\sigma S$ injects new constraint structure from the
accessible-future-volume reservoir. The term $-\lambda\nab\Phi$ in \eqref{eq:ms-v}
drives the inferential current down the constraint gradient (Theorem~\ref{thm:gradient-flow}).
The vorticity term $\kappa_v(\nab\times\vfield)$ sustains the topological structure
of observer identity (Theorem~\ref{thm:vorticity} below). In \eqref{eq:ms-s},
$-\mu\Phi$ encodes Proposition~\ref{prop:contraction}: high constraint density
reduces accessible future volume.

\section{Well-Posedness and Lamphrodyne Smoothing}
\label{sec:wellposed}

\begin{theorem}[Local classical well-posedness]
Under standard regularity and growth conditions, there exists $T^*>0$ such that
the Master System admits a unique classical solution in
$C([0,T^*]; H^2)\times[C([0,T^*]; H^2)]^3\times C([0,T^*]; H^1)$.
\end{theorem}

The proof rewrites the system as $\partial_t u = Lu + N(u)$, applies the abstract
Cauchy theorem for analytic semigroups generated by
$L = (-\kappa_\Phi\Delta, -\kappa_v\Delta, -D_S\Delta)$, and controls the
nonlinearity via the Sobolev embedding $H^1\hookrightarrow L^6$ in $d=3$.

\begin{theorem}[Lamphrodyne smoothing]
\label{thm:smoothing}
If $\sigma\equiv 0$ on $[T,\infty)$, then
$\|\nab S(t)\|_{L^2} \leq \|\nab S(T)\|_{L^2} e^{-\mu_0(t-T)}$,
where $\mu_0 = \min(\mu, D_S\lambda_1)$.
\end{theorem}

\begin{proof}
Multiply \eqref{eq:ms-s} (with $\sigma=0$) by $-\Delta S$ and integrate. Poincar\'e
gives $\tfrac{1}{2}\dot{y} + (D_S\lambda_1+\mu)y \leq 0$ where
$y = \|\nab S\|^2_{L^2}$, yielding the exponential bound.
\end{proof}

\section{Constraint Equilibria}
\label{sec:equilibria}

\begin{theorem}[Constraint equilibrium]
\label{thm:equilibrium}
A stationary solution satisfies $\nab S = 0$, $\nab\Phi = 0$, $\vfield = \mathbf{0}$
if and only if the admissibility landscape contains no preferred future direction.
\end{theorem}

\begin{proof}
Setting all time derivatives to zero: $\lambda\nab\Phi + \nu\vfield = \mathbf{0}$
forces $\vfield = \mathbf{0}$ if $\nab\Phi = \mathbf{0}$; substituting into the
entropy equation gives $\Delta S = 0$; under Neumann conditions the only solution
is constant $S$, hence $\nab S = \mathbf{0}$. Conversely, vanishing gradients imply
no directional preference.
\end{proof}

\section{Vorticity Persistence}
\label{sec:vorticity}

\begin{theorem}[Vorticity persistence]
\label{thm:vorticity}
Let $\omega = \nab\times\vfield$. In the stiff limit with $\nab\cdot\vfield = 0$,
the circulation $\Gamma_C = \oint_C \vfield\cdot\dd\ell$ is conserved along material
loops $C$.
\end{theorem}

\begin{proof}
Kelvin's circulation theorem states $\dd\Gamma_C/\dd t = 0$ for incompressible flow
with conservative body forcing. In the stiff limit, $\nab\cdot\vfield = 0$ (exact
incompressibility) and the forcing $-\lambda\nab\Phi$ is a gradient (hence
conservative). By Stokes' theorem $\Gamma_C = \int_A \omega\cdot\dd A$, vortical
structures persist under advection.
\end{proof}

Theorem~\ref{thm:vorticity} makes the observer-as-vortex interpretation a theorem
rather than a metaphor. A \CLIO\ observer corresponds to a vortical region of the
incompressible constraint flow. Theorem~\ref{thm:vorticity} proves that such a
region is topologically persistent under the Master System dynamics: the circulation
--- and therefore the identity of the observer as a distinct vortex --- is conserved.
Loss of observer identity (agency collapse) requires the forcing to become
non-conservative, which occurs precisely at entropy saturation when $S\to S_{\max}$
and the gradient $\nab\Phi$ becomes dominated by flat regions.

\section{The Stiff Limit and Transverse Flow}
\label{sec:stiff}

In the stiff limit $\kappa_v\to\infty$, the divergence constraint becomes exact.
The Dirac-bracket calculation projects onto transverse flow:
$\{v_i(\mathbf{x}),\pi^j_v(\mathbf{y})\}^* = P^j_i(\mathbf{x}-\mathbf{y})$,
where $P^j_i$ is the transverse projector onto divergence-free modes. The stiff-limit
\RSVP\ is a theory of incompressible, entropy-modulated vortex flow.

% ============================================================
\chapter{The Projection Ladder}
\label{ch:projection}
% ============================================================

\section{The Four Projections}
\label{sec:four}

\begin{definition}[Projection operators]
The canonical projections of $(\Phi, \vfield, S)$, constructed via the kernel
family~\eqref{eq:kernel} with domain-specific choices of $w_i$, are:
\begin{align*}
  \mathcal{P}_{\mathrm{cosmo}} &: w_i = \delta(x-x_i)
    &&\text{(local density, redshift, power spectrum)},\\
  \mathcal{P}_{\mathrm{inst}} &: w_i = m_i(x)
    &&\text{(membership weighting, role observables)},\\
  \mathcal{P}_{\mathrm{obs}} &: w_i = \psi_i(x)
    &&\text{(electrode/BOLD sensitivity kernel)},\\
  \mathcal{P}_{\mathrm{sem}} &: w_i = \phi_i(x)
    &&\text{(concept activation kernel)}.
\end{align*}
\end{definition}

All four projections are instances of the constructive kernel family of
Theorem~\ref{thm:kernels}. The abstract notation of earlier definitions is now
grounded in a concrete construction. The differences among the projections lie
entirely in the choice of kernel $w_i$, not in any separate theoretical apparatus.

Four structural invariants are preserved across all projections: the admissibility
geometry, the constraint structure, the ordering of accessible future volumes
(Theorem~\ref{thm:order}), and causal direction.

\begin{theorem}[Deterministic entropy monotonicity under projection]
For any deterministic projection $\mathcal{P}$,
$H(\mathcal{P}(S)) \leq H(S)$.
\end{theorem}

This follows immediately from Theorem~\ref{thm:det-cg}.

\section{The CLIO Projection}
\label{sec:clio}

\begin{definition}[CLIO descent]
\[
  \mathcal{L}_{\CLIO} = \mathbb{E}[S] + \lambda_{\mathrm{sparse}}\|\Phi\|_0
    - \lambda_{\mathrm{flow}}\|\vfield\|^2.
\]
\end{definition}

\begin{theorem}[Agency requires flow]
If $\lambda_{\mathrm{flow}} > 0$, then $\vfield_t \not\to \mathbf{0}$ along \CLIO\
trajectories unless $\Phi$ is flat.
\end{theorem}

\begin{proof}
$\vfield\to\mathbf{0}$ requires $\nab\Phi\to\mathbf{0}$ (from \eqref{eq:ms-v}),
making $\Phi$ constant. But then the initial contribution
$-\lambda_{\mathrm{flow}}\|\vfield_0\|^2 < 0$ shows $\mathcal{L}$ is strictly
lower than at $\vfield=\mathbf{0}$, contradicting minimality unless $\Phi$ is flat.
\end{proof}

\begin{proposition}[CLIO sparsity / Occam compression]
Among all constraint landscapes generating identical observations, the \CLIO\
objective selects the minimum-description representation.
\end{proposition}

\begin{proof}
For two equivalent fields $\Phi_1, \Phi_2$ producing identical observations, the
entropy and flow terms in $\mathcal{L}_{\CLIO}$ are equal. The optimisation reduces
to $\min\|\Phi\|_0$, selecting the sparsest equivalent representation.
\end{proof}

The neuroimaging prediction: the agency coherence metric $\Gamma(t) =
I(\Phi;\vfield)/H(S)$, computed from the \CLIO-projected fields at the electrode
kernel resolution $\psi_i$, should correlate with EEG microstate duration. This is
an in-principle prediction; the closure relation connecting the field-theoretic
$\Gamma$ to the empirical microstate duration requires the complete constructive
specification of $\psi_i$ for the given electrode array.

\section{TARTAN}
\label{sec:tartan}

\TARTAN\ (Trajectory-Aware Recursive Tiling with Annotated Noise) models memory as
admissibility navigation. Gradient-weighted noise $\eta_s\sim\mathcal{N}(0,c|\nab\Phi|^2)$
concentrates exploration at the steepest constraint boundaries.

\begin{theorem}[Gradient-weighted noise maximises boundary exploration]
The noise variance $c|\nab\Phi|^2$ maximises expected Fisher information near
decision boundaries.
\end{theorem}

\section{Spherepop}
\label{sec:spherepop}

An institution $\mathcal{S} = (I, B, \Sigma)$ with permeability $\Sigma$ and
interior observable $O_I = \int m(x)\Phi(x)\, \dd x$ (from the kernel construction).
The merge condition $\Sigma_i + \Sigma_j > \Theta_{\mathrm{closure}}$ and
$R_t + 2H < P_{\max}$ is a statement about the overlap of membership kernels $m_i$
and the work required to align the constraint landscapes. The institutional merge
prediction: pre-modern merger rates (with $H\approx 0$) should follow the predicted
merge probability. Specifying the permeability $\Sigma$ from the field theory requires
a closure relation for the boundary operator $B$ in terms of the gradient of $\Phi$
at $\partial\Omega_i$.

% ============================================================
\chapter{Emergent Gravity and the Free Energy Principle}
\label{ch:gravity}
% ============================================================

\section{The Jacobson Pathway}
\label{sec:jacobson}

Energy and entropy currents: $J_E\propto\Phi\vfield$, $J_S\propto S\vfield$.
Acceleration $a=\kappa/(2\pi)$ gives Unruh temperature $T=\kappa/(2\pi)$.

\begin{theorem}[Emergence of the Einstein equation]
In the stiff limit, after coarse-graining over a local Rindler horizon and applying
the Clausius relation $\delta Q = T\,\delta S_{\mathrm{horizon}}$ to $J_S$,
\[
  R_{\mu\nu} - \tfrac{1}{2}R\,g_{\mu\nu} + \Lambda g_{\mu\nu} = 8\pi G T_{\mu\nu},
\]
where $g^{\mathrm{eff}}_{\mu\nu} = \mathrm{diag}(-c_s^2+|\vfield|^2,1,1,1)/c_s$.
\end{theorem}

\begin{theorem}[Verlinde limit]
In the quasi-static limit, $m\mathbf{a}\approx g_1\Phi\nab S + g_2 S\nab\Phi$,
Verlinde's entropic force law.
\end{theorem}

\begin{theorem}[Friston-RSVP correspondence]
The \CLIO\ loss is formally equivalent to the variational free energy under
$S\mapsto -\log p$ and $\|\Phi\|_0\mapsto$ model complexity. The correspondence
is an isomorphism of formal structure, not a claim that the physical systems are
identical.
\end{theorem}

The Friston correspondence establishes that active inference is the finite-dimensional,
observer-projected limit of infinite-dimensional \RSVP\ field dynamics. The qualifier
``correspondence'' rather than ``duality'' is used precisely, following the reviewer's
correction.

% ============================================================
\chapter{Cosmology as One Projection}
\label{ch:cosmo}
% ============================================================

\section{The Cosmic Constraint Triad}
\label{sec:triad}

Cosmology is one projection of the admissibility geometry. The allowed cosmos is
the intersection of three deep constraints: geometry $G$ (curvature bounded by CMB
and BAO observations), energy density $E$ (content bounded by nucleosynthesis and
expansion history), and information $I$ (entropy content bounded by the CMB
anisotropy spectrum and the second law). Initial conditions are not fine-tuned in
the traditional sense; they are geometrically selected by the intersection
$G\cap E\cap I$.

\section{Lamphron and Lamphrodyne}
\label{sec:lamphron}

The lamphron process (inward gravitational collapse releasing binding energy into
$\Phi$) and the lamphrodyne process (outward entropic smoothing in void regions)
are twin expressions of the same entropic dynamics at different points in field
configuration space. The cosmological prediction is that redshift, the cosmic web,
and CMB anisotropies are all emergent features of this reorganisation without a
global scale factor.

\section{Entropic Redshift: Parametric Prediction}
\label{sec:redshift}

\begin{definition}[RSVP entropic redshift]
\[
  z_{\RSVP}(t_e, t_0) = \exp\!\left(\alpha_S
    \int_{t_e}^{t_0}\overline{\partial_t S}\, \dd t\right) - 1.
\]
\end{definition}

To make this predictive, a closure relation for $\overline{\partial_t S}(t)$ is
needed. A minimal toy closure assumes exponential decay of the spatial entropy
gradient at rate $H_S$: $\partial_t S = H_S S$, giving
\begin{equation}
  z_{\mathrm{toy}} = \exp(\alpha_S H_S \Delta t) - 1.
  \label{eq:toy-redshift}
\end{equation}
At low redshift this matches the Hubble law with effective $H_0 = \alpha_S H_S$.
Departures from the \LCDM\ prediction at $z\gtrsim 1$ then arise from the
higher-order terms in the entropy equation of state. Equation~\eqref{eq:toy-redshift}
is a parametric toy model, not a finished experimental protocol; its purpose is to
demonstrate that the redshift formula is calculable given a closure relation.

\begin{proposition}[Observable signature]
The entropic redshift formula predicts a departure from \LCDM\ at the $2\sigma$ level,
testable via Type~Ia supernovae, if the entropy production rate $H_S$ and calibration
constant $\alpha_S$ jointly satisfy the relation specified by the Wilks criterion.
The theory is ruled out if $\hat{\alpha}_S < 0$ or if the likelihood ratio exceeds
the Wilks threshold. These conditions constitute the in-principle falsification
criterion; the constructive protocol requires the completed equation of state for
$\overline{\partial_t S}(z)$.
\end{proposition}

% ============================================================
\chapter{The Deck-0 Reservoir: Inexhaustible Refinement}
\label{ch:deck0}
% ============================================================

\section{The Category Error}
\label{sec:category-error}

Lamphrodyne smoothing guarantees gradient decay in the absence of sources. Simulations
require the entropy re-injection $J_{\mathrm{exchange}}$. The energy-battery
interpretation of this injection faces two puzzles: the energy source and the
non-exhaustion. The correct diagnosis is a category error: these are energy questions
about a possibility space.

\section{Inexhaustible Refinement}
\label{sec:inexhaustible}

\begin{theorem}[No terminal refinement]
\label{thm:no-terminal}
Let $\Adm$ be nontrivial. Then there is no terminal partition of $\Adm$ under
admissibility-preserving refinement.
\end{theorem}

\begin{proof}
Assume a terminal partition $\Pi^*$ exists. Choose any cell $C\in\Pi^*$. Since
$\Adm$ is nontrivial, at least two admissible trajectories pass through $C$ with
distinct futures, yielding a strictly finer partition. Contradiction.
\end{proof}

\begin{theorem}[Refinement generates unbounded distinguishability]
\label{thm:unbounded}
Let $\Pi_n$ be a sequence of admissibility-preserving refinements. Then
$|\Pi_{n+1}| > |\Pi_n|$ for all nonterminal partitions, and
$\lim_{n\to\infty}|\Pi_n| = \infty$.
\end{theorem}

\begin{proof}
By Theorem~\ref{thm:no-terminal}, every nonterminal partition $\Pi_n$ admits a
strictly finer partition $\Pi_{n+1}$. Hence $|\Pi_{n+1}| > |\Pi_n|$ at every step.
Since $|\Pi_n|$ is a strictly increasing sequence of positive integers, it is
unbounded.
\end{proof}

Theorem~\ref{thm:unbounded} directly supports the philosophical claim that exploration
creates distinctions. Each act of exploration generates at least one new admissible
trajectory class that did not exist as a distinct possibility before the selection
was made.

\begin{theorem}[Constraint-flow universality class]
\label{thm:universality-class}
Every \RSVP\ projection preserves: admissibility ordering, entropy monotonicity,
gradient-directed flow, and refinement nontermination. Therefore all \RSVP\
projections belong to the same universality class of constraint-governed systems.
\end{theorem}

\begin{proof}
Admissibility ordering is preserved by Theorem~\ref{thm:order}. Entropy monotonicity
follows from Theorem~\ref{thm:det-cg}. Gradient-directed flow is the consequence of
Theorem~\ref{thm:gradient-flow} applied at the projected scale. Refinement
nontermination follows from Theorem~\ref{thm:no-terminal}, since the quotient
admissibility manifold $\Adm/{\sim_{\mathcal{P}}}$ inherits the nontriviality of
$\Adm$. All four properties are intrinsic to the geometry of the admissibility
manifold, hence universal across projections.
\end{proof}

% ============================================================
\chapter{Failure Modes and Structure Formation}
\label{ch:failure}
% ============================================================

\section{The Turing Instability}
\label{sec:turing}

\begin{theorem}[Turing instability threshold]
If $\kappa_v\ll\kappa_\Phi$, a long-wavelength instability emerges when
$\Phi^*>\kappa_v(\kappa_\Phi+D_S)/\kappa_\Phi$.
\end{theorem}

Once the inferential current slows relative to the constraint density, the uniform
admissibility landscape necessarily differentiates. Galaxies, cognitive observers,
and institutions are manifestations of this instability at different projection
scales.

\section{Unified Failure Modes}
\label{sec:failures}

Three canonical failure modes appear at every level of the projection ladder.

\textit{Entropy saturation.} $S\to S_{\max}$ uniformly. Accessible future volume
exhausted. No new distinctions possible. In cosmology: heat death. In cognition:
representational collapse. In institutions: homogenisation.

\textit{Gradient flattening.} $|\nab S|\to 0$ without saturation. Directional bias
lost. The vanishing-gradient analogue.

\textit{Vorticity decay.} $\curlop\vfield\to\mathbf{0}$. Persistence of observer
identity fails. By Theorem~\ref{thm:vorticity}, this requires non-conservative
forcing --- precisely the condition that arises at entropy saturation. Hence vorticity
decay is the first warning sign; entropy saturation is the terminal state.

% ============================================================
\chapter{Open Problems}
\label{ch:open}
% ============================================================

The research frontier is defined by the following open problems, ordered by impact.

\textit{Open Problem~1 (General Unruh theorem).} Derive the general result
$T_{\mathrm{eff}} \propto \dd(\log\Vac_{\mathrm{hidden}})/\dd t$ from the \RSVP\
field equations, moving the Unruh principle from analogy to derivation.

\textit{Open Problem~2 (Constructive projection functors).} Specify the kernel
functions $w_i$ for each domain from the Master System dynamics --- that is, derive
the electrode sensitivity profiles, membership weightings, and cosmological response
functions from the field theory rather than importing them from domain expertise.

\textit{Open Problem~3 (Entropy equation of state).} Derive $\overline{\partial_t S}(z)$
as a function of cosmic time or redshift from the \RSVP\ field equations. The toy
closure~\eqref{eq:toy-redshift} provides the structure; the complete derivation
would constitute a finished experimental protocol for the entropic redshift
prediction.

\textit{Open Problem~4 (Measure on trajectory space).} Identify the appropriate
measure $\mathcal{D}\gamma$ on the admissibility manifold. In the Minkowski vacuum
it is the Euclidean path integral measure; in a curved spacetime with dynamical
entropy the appropriate measure is unknown.

\textit{Open Problem~5 (Quantisation).} Promote the \RSVP\ fields to operators on
a Hilbert space. The constrained Hamiltonian structure in the stiff limit provides
the starting point; the nonlinearity of the entropy equation poses significant
challenges.

% ============================================================
\chapter{Conclusions}
\label{ch:conclusions}
% ============================================================

The deepest primitive of the \RSVP\ framework is not matter, energy, space,
information, or entropy in the traditional senses. It is the topology of what can
still happen.

That sentence opened this monograph. It can now be given a precise meaning in terms
of the objects developed here. The topology of what can still happen is the
admissibility manifold $\Adm$. Its shape is indexed by the constraint density $\Phi$.
Its measure is the accessible future volume $\Vac$, uniquely logarithmic by
Proposition~\ref{prop:log-unique}. Its navigation is driven by the inferential
current $\vfield$. Its persistent structures are vortices, proven topologically
stable by Theorem~\ref{thm:vorticity}. Its failure modes are entropy saturation,
gradient flattening, and vorticity decay. Its inexhaustibility is proven by
Theorem~\ref{thm:no-terminal} and Theorem~\ref{thm:unbounded}. Its universality
across substrates is established by Theorem~\ref{thm:inversion}, and its universality
class is characterised by Theorem~\ref{thm:universality-class}.

The historical trajectory of the framework:
Newton started with forces and implied geometry.
Einstein started with geometry and implied thermodynamics.
Friston starts with thermodynamics and implies inference.
\RSVP\ starts with physical relaxation, discovers thermodynamic structure, discovers
accessibility structure, and arrives at constraint geometry.

The cosmological interpretation may survive or fail observational testing. The
entropic redshift hypothesis will be confirmed or rejected by forthcoming supernova
surveys. But the deeper contribution is independent of those outcomes: a framework
in which the admissibility manifold is the primitive object, projection is a
first-class mathematical operation, entropy is derived (not postulated), and observer
identity is a theorem of incompressible vortex flow rather than a metaphor.

Reality consists of systems navigating constraint landscapes, exploring admissible
futures, and generating new distinctions. The cosmological theory is one application
of that principle. The principle itself is what the theory has become.

\appendix

\chapter{Euler-Lagrange Equations}
\label{app:el}

\begin{align*}
\ddot{\Phi} - c_\Phi^2\Delta\Phi + U'_\Phi &=
  \kappa_v\alpha_\Phi(\nab\cdot\vfield - \alpha_\Phi\Phi - \alpha_S S)
  + g_1\vfield\cdot\nab S + g_2 S(\nab\cdot\vfield) + g_3 S,\\
\ddot{\vfield} + c_v^2\nab\times(\nab\times\vfield) &=
  \kappa_v(\nab\cdot\vfield - \alpha_\Phi\Phi - \alpha_S S)\nab(\nab\cdot\vfield)
  + g_1\Phi\nab S + g_2 S\nab\Phi,\\
\ddot{S} - c_S^2\Delta S + U'_S &=
  \kappa_v\alpha_S(\nab\cdot\vfield - \alpha_\Phi\Phi - \alpha_S S)
  + g_2\vfield\cdot\nab\Phi + g_1\nab\cdot(\Phi\vfield) + g_3\Phi.
\end{align*}

\chapter{Conservation Laws}
\label{app:conservation}

Total constraint density $\int_\Omega\Phi\, \dd x$ is conserved. Energy-entropy
identities:
\begin{align*}
\tfrac{\dd E_\Phi}{\dd t} &= -\lambda\textstyle\int\Phi(\nab\cdot\vfield)
  +\sigma\int\Phi S - \mu\int\Phi^2,\\
\tfrac{\dd E_v}{\dd t} &= -\lambda\int\vfield\cdot\nab\Phi - \nu\int|\vfield|^2,\\
\tfrac{\dd E_S}{\dd t} &= -\sigma\int S\Phi - \eta\int S^2.
\end{align*}
The sum is non-positive under the positivity conditions on coupling constants.

\chapter{RSVP Labs: Reduction Pathways}
\label{app:labs}

The \RSVP\ Labs 1--40 are a suite of numerical simulations each implementing one
reduction of the Master System. They are grouped as follows. Scalar-only
($\vfield=\mathbf{0}$): $\partial_t\Phi = \sigma S - U'_\Phi$ (Labs 3,7,9,15,21,32).
Vector-only ($\Phi=0$): $\partial_t\vfield = -\nu\vfield + \kappa(\nab\times\vfield)$
(Labs 1,27,38). Entropy-reservoir coupling, Deck-0 dynamics:
$\partial_t S = -\mu\Phi + \chi|\vfield|^2$ (Labs 6,14,24,30). Reaction-diffusion
with morphogen activator-inhibitor structure (Labs 19,38). Phase-space reduction
to attractor with Kuramoto synchronisation (Labs 13,20,23,25,29,39). Observer
projection via electrode kernels (Labs 22,35,37,40).

\backmatter

\chapter*{Acknowledgements}
\addcontentsline{toc}{chapter}{Acknowledgements}

This work was developed as an independent research programme. Conceptual debts are
owed to Ted Jacobson, Erik Verlinde, Karl Friston, William Unruh, Amal Raychaudhuri,
and the traditions of analogue gravity, nonequilibrium thermodynamics, and effective
field theory for continuous media. The internal and external peer review that
motivated the major structural changes in this version --- the derivation of the
logarithm, the constructive kernel family, the split entropy monotonicity theorems,
the vorticity persistence theorem, the unbounded distinguishability theorem, the
universality class theorem, the geometry of projection chapter, and the toy cosmological
closure --- were essential. The author bears sole responsibility for errors.

\chapter*{References}
\addcontentsline{toc}{chapter}{References}

\begin{thebibliography}{30}

\bibitem{jacobson1995}
Jacobson, T. (1995). Thermodynamics of spacetime: The Einstein equation of state.
\textit{Physical Review Letters}, 75(7), 1260--1263.

\bibitem{verlinde2011}
Verlinde, E. (2011). On the origin of gravity and the laws of Newton.
\textit{Journal of High Energy Physics}, 2011(4), 029.

\bibitem{unruh1976}
Unruh, W.~G. (1976). Notes on black hole evaporation.
\textit{Physical Review D}, 14(4), 870--892.

\bibitem{hawking1975}
Hawking, S.~W. (1975). Particle creation by black holes.
\textit{Communications in Mathematical Physics}, 43(3), 199--220.

\bibitem{raychaudhuri1955}
Raychaudhuri, A.~K. (1955). Relativistic cosmology I.
\textit{Physical Review}, 98(4), 1123--1126.

\bibitem{padmanabhan2005}
Padmanabhan, T. (2005). Gravity and the thermodynamics of horizons.
\textit{Physics Reports}, 406, 49--125.

\bibitem{barcelo2005}
Barcel\'o, C., Liberati, S., \& Visser, M. (2005). Analogue gravity.
\textit{Living Reviews in Relativity}, 8, 12.

\bibitem{volovik2003}
Volovik, G.~E. (2003). \textit{The Universe in a Helium Droplet}. Oxford UP.

\bibitem{weinberg1995}
Weinberg, S. (1995). \textit{The Quantum Theory of Fields}, Vol.~I. Cambridge UP.

\bibitem{nicolis2013}
Nicolis, A., \& Penco, R. (2013). Mutual interactions of phonons, rotons, and
gravity. \textit{Physical Review D}, 87(12), 123007.

\bibitem{andersson2007}
Andersson, N., \& Comer, G.~L. (2007). Relativistic fluid dynamics.
\textit{Living Reviews in Relativity}, 10, 1.

\bibitem{onsager1931}
Onsager, L. (1931). Reciprocal relations in irreversible processes I.
\textit{Physical Review}, 37(4), 405--426.

\bibitem{prigogine1967}
Prigogine, I. (1967). \textit{Introduction to Thermodynamics of Irreversible
Processes}. Wiley.

\bibitem{pazy1983}
Pazy, A. (1983). \textit{Semigroups of Linear Operators}. Springer.

\bibitem{lions1969}
Lions, J.-L. (1969). \textit{Quelques m\'ethodes de r\'esolution}. Dunod.

\bibitem{dirac1964}
Dirac, P.~A.~M. (1964). \textit{Lectures on Quantum Mechanics}. Yeshiva.

\bibitem{henneaux1992}
Henneaux, M., \& Teitelboim, C. (1992). \textit{Quantization of Gauge Systems}.
Princeton UP.

\bibitem{bekenstein1973}
Bekenstein, J.~D. (1973). Black holes and entropy.
\textit{Physical Review D}, 7(8), 2333--2346.

\bibitem{sakharov1968}
Sakharov, A.~D. (1968). Vacuum quantum fluctuations and the theory of gravitation.
\textit{Soviet Physics Doklady}, 12, 1040.

\bibitem{friston2021}
Friston, K.~J. (2021). The active inference framework.
\textit{Nature Reviews Neuroscience}, 22(3), 167--187.

\bibitem{barandes2024}
Barandes, J.~A. (2024). New prospects for a causally local formulation of quantum
theory. \textit{arXiv:2402.16935}.

\bibitem{planck2020}
Planck Collaboration (2020). Planck 2018 results VI.
\textit{Astronomy \& Astrophysics}, 641, A6.

\bibitem{riess2022}
Riess, A.~G., et al. (2022). A comprehensive measurement of $H_0$.
\textit{Astrophysical Journal Letters}, 934(1), L7.

\bibitem{vopson2020}
Vopson, M.~M. (2020). The information catastrophe.
\textit{AIP Advances}, 10(8), 085014.

\bibitem{bianconi2025}
Bianconi, G. (2025). Gravity from entropy.
\textit{Physical Review D}, 111, 066001.

\bibitem{carroll2016}
Carroll, S.~M., \& Remmen, G.~N. (2016). What is the entropy in entropic gravity?
\textit{Physical Review D}, 93(12), 124052.

\bibitem{visser2011}
Visser, M. (2011). Conservative entropic forces.
\textit{Journal of High Energy Physics}, 2011(10), 140.

\bibitem{cheung2008}
Cheung, C., et al. (2008). The EFT of inflation.
\textit{Journal of High Energy Physics}, 2008(3), 014.

\bibitem{endlich2013}
Endlich, S., Nicolis, A., \& Penco, R. (2013). UV completion without symmetry
restoration. \textit{Physical Review D}, 88(10), 105001.

\bibitem{nesterov2004}
Nesterov, Y. (2004). \textit{Introductory Lectures on Convex Optimization}.
Springer.

\end{thebibliography}

\end{document}
